% ECONOMICS_PROBLEM: {"id":"C61-73","title":"Random maturity and filtering in joint control and stopping","jel_code":"C61","source_id":"OP-0570"}
% FIRST_POSTED: 2026-10-10
% ENTRY_KIND: research_attempt
% ATTEMPT_STATUS: solved
% !TeX program = xelatex
\documentclass[11pt,a4paper]{article}
\usepackage[margin=25mm]{geometry}
\usepackage{fontspec}
\setmainfont{Latin Modern Roman}
\usepackage{amsmath,amssymb,amsthm}
\usepackage{mathrsfs}
\usepackage{hyperref}
\hypersetup{colorlinks=true,linkcolor=blue,citecolor=blue,urlcolor=blue}
\setlength{\emergencystretch}{3em}
\newtheorem{theorem}{Theorem}[section]
\newtheorem{lemma}[theorem]{Lemma}
\newtheorem{proposition}[theorem]{Proposition}
\newtheorem{corollary}[theorem]{Corollary}
\theoremstyle{definition}
\newtheorem{definition}[theorem]{Definition}
\theoremstyle{remark}
\newtheorem{remark}[theorem]{Remark}
\newcommand{\R}{\mathbb R}
\newcommand{\E}{\mathbb E}
\newcommand{\Qq}{\mathbb Q}
\newcommand{\Ptwo}{\mathcal P_2(\R)}
\newcommand{\CW}{C_{\mathsf w}}
\newcommand{\wt}{\mathsf w}
\newcommand{\cQ}{\mathcal Q}
\newcommand{\cT}{\mathcal T}
\newcommand{\cH}{\mathcal H}
\newcommand{\ind}{\mathbf 1}
\newcommand{\sech}{\operatorname{sech}}
\DeclareMathOperator*{\argmin}{arg\,min}
\title{Random Maturity, Survival Filtering, and Joint Control and Stopping\\
\large A posterior obstacle characterization for C61-73}
\author{Ji Ho Bae\thanks{Prepared with GPT Codex assistance.}}
\date{10 October 2026}
\begin{document}
\maketitle
\begin{center}\small Complete hand proof with internal adversarial review\end{center}

\begin{abstract}
For the full hidden-diffusion model specified in C61-73, we characterize
the value by constant-control killed Bayes kernels on the space of
posterior probability measures. A killed second-moment weight makes
every grid Bellman operator a strict contraction. Its exact posterior
policies converge in value to the original continuous-time problem,
without convexifying the control cost or identifying observation and
innovation filtrations. The value is the greatest continuous weighted
cost-submean minorant of the stopping obstacle and is also the unique
leavable-semigroup fixed point. These are verified mild meanings of the
measure-valued obstacle equation and require no classical solution or
viscosity comparison assumption. We give an exact-optimality certificate
and pure posterior grid rules arbitrarily close to optimal. A nonconstant
bounded hazard admits an explicit one-dimensional survival filter and
a stopping threshold, while the general characterization retains all
coefficients allowed in the canonical problem.
\end{abstract}

\section{The specified problem and the form of the answer}

Let \(U\subset\R\) be a compact interval. In the original model,
\[
 dX_t=u_t\mu(X_t)\,dt+\sigma(X_t)\,dW_t,\qquad
 dO_t=h(X_t)\,dt+dV_t,
\]
where \(W,V\) are independent Brownian motions. The functions
\(\mu,\sigma,h\) are globally Lipschitz, \(\mu,h\) are bounded, and
\(\sigma\) is bounded and bounded away from zero. Let \(a\) be bounded
Lipschitz, \(a\geq a_0>0\), and \(E_0\) an independent unit exponential
variable. The initial law \(\pi\in\Ptwo\) is independent of these
noises and \(E_0\). Set
\[
 A_t=\int_0^t a(X_s)\,ds,\quad
 \eta=\inf\{t:A_t\geq E_0\},\quad N_t=\ind_{\{\eta\leq t\}}.
\]
The information filtration \(\mathcal G\) is the usual augmentation
of that generated by \((O,N)\). Controls are \(U\)-valued progressive
rules for this information; stopping times are \(\mathcal G\)-stopping
times. As usual, stochastic-integral admissibility uses a progressively
measurable version.

The nonnegative costs \(k,\psi\) are continuous, \(k(x)\leq C(1+x^2)\),
and \(c>0\). Write \(\ell(u)=\psi(u)+c\). We characterize
\begin{equation}\label{eq:original}
 V(\pi)=\inf_{u,\tau}\E_\pi\left[
 k(X_{\tau\wedge\eta})+\int_0^{\tau\wedge\eta}\ell(u_s)\,ds\right].
\end{equation}
Voluntary stopping at infinity is allowed; maturity is almost surely
finite because \(a\geq a_0\).

The canonical convention permits \(\varepsilon\)-optimal rules when
attainment is not assured. Our general result gives the full value and
such pure admissible rules for every \(\varepsilon>0\), and an
exact-optimality certificate. We do not assume convexity of \(\psi\),
a finite-dimensional filter, an independent hazard, a classical value
function, or an infinite-dimensional viscosity comparison theorem.
The primary obstacle answer is a greatest cost-submean minorant defined
from known constant-control kernels. An equivalent semigroup fixed
point is proved. The result is therefore not only a verification
statement conditional on existence of a smooth value function.

\section{Pre-default policies and a fixed observation reference}

\begin{lemma}[Raw pre-default reduction]\label{lem:raw}
Up to changes irrelevant to \eqref{eq:original}, every original policy
has a raw observation functional \(u^0(O)\) before maturity and an
observation stopping functional \(\tau^0(O)\) such that
\[
 u_t=u_t^0(O)\quad(t<\eta),\qquad
 \tau\wedge\eta=\tau^0(O)\wedge\eta
 \quad\hbox{almost surely}.
\]
These functionals of a continuous observation path are nonanticipating.
\end{lemma}
\begin{proof}
For a raw canonical rule on paths \((o,n)\), substitute the identically
zero path for \(n\). Before the actual default jump, the observed path
restrictions agree with those of the zero-default path. Nonanticipation
gives the assertions. Measurability is first checked for predictable
rectangles or progressive cylinder functions and then extended by a
monotone-class argument.

For a stopping time, use the canonical rational-time events specifying
its dyadic upper approximations; their decreasing limit defines the
observation stopping functional. A stop exactly at default gives the
same minimum with \(\eta\). Usual completions are removed by raw versions
in these countable approximations. On every finite horizon the reference
observation law below and the surviving observation law have strictly
positive densities relative to Wiener measure. Thus exceptional sets
can be discarded without changing any pre-default cost.
\end{proof}

It is important not to integrate \(E_0\) after conditioning on the entire
original signal path: that signal may depend on default through the
post-default control. The following stopped construction avoids this.

\begin{lemma}[Stopped reverse Girsanov representation]\label{lem:reverse}
Every original admissible policy has the pre-default cost representation
obtained from an exogenous Brownian observation \(Y\), independent of the
signal Brownian motion and of \((X_0,E_0)\), with control \(u^0(Y)\).
\end{lemma}
\begin{proof}
Fix \(T<\infty\). On the original space put
\[
 \widetilde Y_t=V_t+\int_0^{t\wedge\eta}h(X_s)\,ds,
\quad
 \Lambda_T=\exp\left\{-\int_0^{T\wedge\eta}h(X_s)\,dV_s
              -\frac12\int_0^{T\wedge\eta}h(X_s)^2\,ds\right\}.
\]
Boundedness of \(h\) makes the exponential a true martingale, also
conditionally on the initial sigma field containing \((X_0,E_0)\).
Under the measure with density \(\Lambda_T\), the pair
\((W,\widetilde Y)\) is jointly Brownian in the full filtration, by
Girsanov's formula and its identity quadratic-covariation matrix.
The law of \((X_0,E_0)\) is unchanged, and this pair is independent of
the Brownian motions.

Before \(\eta\), \(\widetilde Y=O\) and \(u=u^0(\widetilde Y)\).
On this reference space construct, throughout \([0,T]\),
\[
 d\overline X_t=u_t^0(\widetilde Y)\mu(\overline X_t)\,dt
                   +\sigma(\overline X_t)\,dW_t,\qquad
 \overline X_0=X_0.
\]
Uniform Lipschitz coefficients give a unique strong solution for this
exogenous observation control. Stopped pathwise uniqueness gives
\(\overline X=X\) on \([0,\eta\wedge T]\). Thus the hazard crossing
computed from \(\overline X\) agrees with \(\eta\) before \(T\), and
its minimum with \(T\) is \(\eta\wedge T\). This reference extension
on \([0,T]\), unlike the original post-default signal, is independent
of \(E_0\).

At \(\theta=\tau^0\wedge\eta\wedge T\), the inverse density is the
unstopped reference observation likelihood built from \(\overline X\)
evaluated at \(\theta\). For a running term on \(\{s<\theta\}\), use
its likelihood at \(s\). We may now integrate the independent exponential
variable against this reference extension. No independence of \(E_0\)
and the original post-maturity signal has been assumed.
\end{proof}

Henceforth let \(\Qq\) carry independent \(X_0\sim\pi,W,Y\), with \(W,Y\)
Brownian. A reference policy is any \(Y\)-progressive \(U\)-valued process
and any \(Y\)-stopping time. Solve
\[
 dX_t=u_t\mu(X_t)\,dt+\sigma(X_t)\,dW_t
\]
and define
\begin{equation}\label{eq:likelihood}
 L_t=\exp\left\{\int_0^t h(X_s)\,dY_s
                     -\frac12\int_0^t h(X_s)^2\,ds\right\},
 \qquad Z_t=L_te^{-A_t}.
\end{equation}
The likelihood is a true martingale on every finite horizon, and
\begin{equation}\label{eq:Lmoment}
 \E_\Qq[L_t^p\mid X_0]\leq
 \exp\left\{\frac{p(p-1)}2\|h\|_\infty^2t\right\}\quad(p>1).
\end{equation}
The same bound holds conditionally on coupled initial variables
independent of the driving noises. The unnormalized and normalized
survival filters are
\begin{equation}\label{eq:rho}
 \rho_t(\varphi)=\E_\Qq[Z_t\varphi(X_t)\mid\mathcal F_t^Y],
 \qquad \pi_t=\rho_t/\rho_t(1).
\end{equation}
Their finite-time versions, including posterior values at stopping
times, are justified below.

\begin{proposition}[Separated cost and maturity tilt]\label{prop:cost}
The reference criterion is
\begin{equation}\label{eq:separated}
 J_\pi(u,\tau)=
 \E_\Qq\left[\rho_\tau(k)+
       \int_0^\tau\{\rho_s(ak)+\ell(u_s)\rho_s(1)\}\,ds\right],
\end{equation}
with zero voluntary terminal term on \(\{\tau=\infty\}\). It represents
every original admissible cost as above. At the maturity jump,
\begin{equation}\label{eq:tilt}
 \E[k(X_\eta)\mid\mathcal G_\eta]
       =\frac{\pi_{\eta-}(ak)}{\pi_{\eta-}(a)}.
\end{equation}
\end{proposition}
\begin{proof}
In the independent reference extension integrate \(E_0\) first.
Voluntary termination contributes \(e^{-A_\tau}k(X_\tau)\), maturity
contributes \(\int_0^\tau e^{-A_s}a(X_s)k(X_s)\,ds\), and survival
multiplies the running cost by \(e^{-A_s}\). Apply the likelihood
and conditional expectation to obtain \eqref{eq:separated}.
Nonnegative truncations establish this identity before integrability
is invoked.

The conditional compensator of the maturity jump and its state mark is
\(\ind_{\{s<\eta\}}\pi_s(dx)a(x)\,ds\).
Its total intensity is \(\pi_s(a)\); normalization gives
\eqref{eq:tilt}. First use bounded truncations of \(k\), then the
quadratic bound in the next section.
\end{proof}

Independent randomization cannot improve an infimum: condition on its
seed and average costs of pure rules. We use reference controls for
value arguments, not to assert that every measurable continuous-time
feedback has a strong realization on a prescribed original Brownian
space. Grid rules will be strongly realized, and their density proves
that the reference enlargement has no effect on the value.

\section{A killed second-moment weight}

Put
\[
 B=\max_{u\in U}|u|\,\|\mu\|_\infty,\quad S=\|\sigma\|_\infty,\quad
 \gamma=a_0/2,\quad R\geq\max\{1,B^2/\gamma^2+S^2/\gamma\},
\]
and \(w(x)=R+x^2,\ \wt(\pi)=\pi(w)\). For
\(\mathcal L^u\varphi=u\mu\varphi'+\sigma^2\varphi''/2\),
Young's inequality gives
\begin{equation}\label{eq:Lyapunov}
 \mathcal L^u w\leq\gamma x^2+B^2/\gamma+S^2\leq\gamma w,
 \qquad(\mathcal L^u-a)w\leq-\gamma w.
\end{equation}

\begin{lemma}[Killed weight and uniform tails]\label{lem:weight}
For all reference policies,
\begin{equation}\label{eq:mass}
 \E_\Qq[\rho_t(w)]\leq e^{-\gamma t}\wt(\pi).
\end{equation}
The conditional optional version holds after any observation history.
Choose \(C_k,C_r<\infty\) so that \(k\leq C_kw\) and
\(ak+\ell(u)\leq C_rw\) for \(u\in U\).
Every policy has finite cost at most
\((C_k+C_r/\gamma)\wt(\pi)\). Replacing \(\tau\) by \(\tau\wedge T\)
changes its cost in absolute value by at most
\begin{equation}\label{eq:tail}
 (2C_k+C_r/\gamma)e^{-\gamma T}\wt(\pi).
\end{equation}
\end{lemma}
\begin{proof}
Itô's formula with localization for \(e^{-A_t}w(X_t)\), under the
physical law of the reference extension, and \eqref{eq:Lyapunov}
give the assertion. Bounded-coefficient moment bounds and
nonnegativity remove localization. Projection equivalently makes
\(e^{\gamma t}\rho_t(w)\) a nonnegative supermartingale. This gives
the conditional optional estimates as well.

Apply these estimates to the voluntary terminal term and integrate
\eqref{eq:mass} for running and maturity costs. On \(\{\tau>T\}\)
bound each of the old and new voluntary terminal terms by
\(C_ke^{-\gamma T}\wt(\pi)\); the discarded integral is at most
\((C_r/\gamma)e^{-\gamma T}\wt(\pi)\). This proves \eqref{eq:tail},
uniformly in the policy.
\end{proof}

\section{Bayes kernels, continuity, and measurable updates}

For a constant \(u\) on \([0,t]\), define
\begin{align}
 \cQ_t^u f(\pi)&=\E_\Qq[\rho_t(1)f(\pi_t)],\label{eq:Q}\\
 g_t(\pi,u)&=\E_\Qq\int_0^t
        [\rho_s(ak)+\ell(u)\rho_s(1)]\,ds,\qquad K(\pi)=\pi(k).
        \label{eq:g}
\end{align}
The kernel is substochastic: its total mass is survival probability.
In particular,
\begin{equation}\label{eq:Qweight}
 \cQ_t^u\wt(\pi)\leq e^{-\gamma t}\wt(\pi).
\end{equation}
Let \(\CW\) be continuous functions on \(\Ptwo\), with its
Wasserstein-\(2\) topology, with finite Banach norm
\[
 \|f\|_{\wt}=\sup_{\pi\in\Ptwo}|f(\pi)|/\wt(\pi).
\]

\begin{lemma}[Weighted Bayes continuity]\label{lem:Feller}
For \(f\in\CW\), both \((\pi,u)\mapsto\cQ_t^u f(\pi)\) and
\((\pi,u)\mapsto g_t(\pi,u)\) are continuous, with the indicated
weighted growth. The same terminal-posterior convergence argument
applies when reference controls converge in \(L^2(dt\,d\Qq)\).
\end{lemma}
\begin{proof}
For \(\pi_n\to\pi\) in Wasserstein-\(2\), couple initial variables
independently of common \(W,Y\) with
\(\E|X_0^n-X_0|^2\to0\). If \(u_n\to u\), the Lipschitz estimate gives
\[
 \E_\Qq\sup_{s\leq t}|X_s^n-X_s|^2\longrightarrow0.
\]
For progressive controls the same estimate adds a constant times
\(\E\int_0^t|u_s^n-u_s|^2ds\), so the stated extension holds.
Bounded Lipschitz \(h,a\) imply convergence of \(Z_t^n\) in probability.
Uniform likelihood moments of all fixed orders imply likelihood
convergence in every fixed finite \(L^p\), by using a higher moment
for uniform integrability.

The quadratic step uses only a second initial moment. Write
\(X_s^n=X_0^n+D_s^n\). Bounded drift and diffusion give uniform
moments of every fixed order of
\(\sup_{s\leq t}|D_s^n|\); their products with likelihoods are uniformly
integrable by Hölder's inequality. For the initial term use
\[
 \E_\Qq[L_s^n\mid X_0^n,X_0]=1
\]
and the uniform conditional \(p\)-moment bound. Split first at
\(|X_0^n|^2>M\), whose weighted expectation is just the ordinary initial
second-moment tail, uniformly small by the \(L^2\) coupling.
On its complement split at \(L_s^n>N\), controlled uniformly by its
higher moment. Consequently
\(\{Z_s^n(R+|X_s^n|^2):n\geq1,\ s\leq t\}\) is uniformly integrable.
The same proof works at bounded observation stopping times.

Conditional expectation contracts \(L^1\). Thus
\(\rho_t^n(\varphi)\to\rho_t(\varphi)\) in \(L^1\) for bounded continuous
\(\varphi\), and for \(\varphi=x^2\) by the preceding product bound.
A countable bounded-Lipschitz determining family, the second moment,
and the positive normalizing mass yield
\(\pi_t^n\to\pi_t\) in Wasserstein-\(2\) in probability.
Hence \(\rho_t^n(1)f(\pi_t^n)\) converges in probability.
Its absolute value is bounded by
\(\|f\|_{\wt}\rho_t^n(w)\), a uniformly integrable family:
conditional expectation preserves uniform integrability by the
de la Vallée-Poussin criterion and Jensen's inequality.
Expectations converge. Integrating the same estimates proves the
assertion for \(g_t\).
\end{proof}

\begin{lemma}[Continuous posterior versions]\label{lem:versions}
The filters have versions for which \(\rho_t(1)>0\) at every finite
time and \(t\mapsto\pi_t\) is continuous in Wasserstein-\(2\).
\end{lemma}
\begin{proof}
For smooth compactly supported tests, the conditional Itô calculation
gives
\[
 d\rho_t(\varphi)=
 \rho_t(\mathcal L^{u_t}\varphi-a\varphi)\,dt+
 \rho_t(h\varphi)\,dY_t,
\]
so these coordinates have continuous versions. On \([0,T]\),
\[
 H=\sup_{s\leq T}L_s\left(1+\sup_{s\leq T}|X_s|^2\right)
\]
is integrable. For its initial-state square term condition on \(X_0\)
and use the conditional likelihood supremum bound; for its increment
terms use higher moments of bounded-coefficient increments and
likelihoods. The tails
\[
 H_R=\sup_{s\leq T}L_s\sup_{s\leq T}|X_s|^2
                 \ind_{\{\sup_{s\leq T}|X_s|>R\}}
\]
have expectations tending to zero. A smooth truncation of \(x^2\)
has conditional-projection error bounded by a constant times
\(\E_\Qq[H_R\mid\mathcal F_t^Y]\). Doob's \(L^1\) maximal inequality
makes this error tend to zero uniformly in time in probability.
A subsequence converges uniformly almost surely, giving a continuous
second-moment coordinate.

The mass solves
\[
 d\rho_t(1)=\rho_t(1)
       [-\pi_t(a)\,dt+\pi_t(h)\,dY_t].
\]
Localization before a possible zero, with bounded logarithmic
coefficients, expresses it as a strictly positive exponential on
every finite interval; it cannot reach zero. Normalize the countable
continuous weak-determining coordinates and the continuous second
moment. Together they imply Wasserstein-\(2\) continuity.
\end{proof}

\begin{lemma}[Jointly Borel sample update]\label{lem:update}
There is a jointly Borel map
\(\Phi_t(\pi,u,y)\in\Ptwo\), with
\(y\in C([0,t],\R)\), giving the posterior after a constant-control
interval from its initial law, action, and full observation increment.
\end{lemma}
\begin{proof}
Represent the initial law by its quantile applied to an independent
uniform variable. The Lipschitz SDE's Picard construction gives a
jointly measurable state solution in this variable, the action, and
the noises. Hence the finite measure
\[
 \E_\Qq[Z_t\ind_{\{X_t\in dx\}}\ind_{\{Y\in C\}}]
\]
depends measurably on \((\pi,u)\) for every Borel observation event \(C\).

Take increasing finite partitions of \(C([0,t],\R)\) generating its
Borel sigma field, using rational-time evaluations and truncated
dyadic intervals. On a positive-Wiener-probability cell divide this
measure by the cell probability and normalize. This gives jointly
Borel finite-partition \(\Ptwo\)-valued conditional measures.
Martingale convergence for \(1,x^2\) and a countable bounded-Lipschitz
family gives Wasserstein-\(2\) convergence for Wiener-almost every
path, for each parameter. The set where the measures are Cauchy and
their unnormalized masses have positive limit is Borel; take their
limit there, and \(\delta_0\) elsewhere. The exceptional set is null
also for the surviving observation law, with density \(\rho_t(1)\).
This remains true for past-measurable initial posteriors and actions:
conditional on the past, the fresh reference observation segment has
Wiener law and is independent of those parameters. Hence the map is
valid for feedback and concatenation.
\end{proof}

\section{Contractive Bellman equations and exact grid policies}

For \(d>0\) set
\begin{equation}\label{eq:T}
 \cT_df(\pi)=\min\left\{K(\pi),\
       \min_{u\in U}[g_d(\pi,u)+\cQ_d^u f(\pi)]\right\}.
\end{equation}

\begin{theorem}[Exact grid problem]\label{thm:grid}
The operator \(\cT_d\) maps \(\CW\) to itself and is a contraction
of constant \(q_d=e^{-\gamma d}<1\). Its unique fixed point \(v_d\)
is the exact value for controls held constant between grid dates
and voluntary stopping on the grid. Moreover \(0\leq v_d\leq K\) and
\begin{equation}\label{eq:griderror}
 \|\cT_d^NK-v_d\|_{\wt}\leq C_kq_d^N.
\end{equation}
A pure posterior feedback attains \(v_d\).
\end{theorem}
\begin{proof}
Lemma~\ref{lem:Feller} and compactness of \(U\) give continuity of
the inner minimum. Equation~\eqref{eq:Qweight} gives
\(\|\cT_df-\cT_dg\|_{\wt}\leq q_d\|f-g\|_{\wt}\).
Geometric convergence of its iterates proves the fixed-point
assertions. Monotonicity and the stopping option give the bounds.

Let \(C_d(\pi)=\min_u[g_d(\pi,u)+\cQ_d^uv_d(\pi)]\).
Stop when \(K(\pi)\leq C_d(\pi)\); otherwise choose the smallest
minimizer and hold it to the next grid date or maturity.
This selector is Borel: the event that it is less than \(r\) is a
countable union of events equating the minimum on a compact initial
subinterval of \(U\) ending below \(r\) with the full minimum.
Both minimum functions are continuous.

Construct this policy strongly from the original \(W,V\), one
interval at a time. The action is known at the interval's beginning;
the state and observation SDEs are Lipschitz there.
Lemma~\ref{lem:update} supplies the Borel update at its end.
Every observation-history grid policy has the same strong construction.

Conditional on the observed past and survival, the signal has the
current posterior law, future Brownian increments are fresh, and
the residual exponential variable is memoryless. Thus \(g_d,\cQ_d^u\)
are the exact conditional next-step cost and posterior law. Past
information beyond the posterior is only additional randomization,
which cannot improve an infimum. This proves grid dynamic programming.

Iterate the fixed-point equality along the displayed policy.
The cost through \(N\) intervals plus its surviving \(v_d\) terminal
term equals \(v_d\). The terminal term is bounded by
\(C_kq_d^N\wt(\pi)\), even if voluntary stopping never occurs.
Letting \(N\to\infty\) proves attainment. The Bellman inequality
gives the lower bound for every history-dependent grid policy.
\end{proof}

\section{Grid density without relaxing the control cost}

\begin{lemma}[Approximation of each reference policy]\label{lem:density}
Every reference policy on a bounded horizon can be approximated in
cost by observation-history policies constant on nested dyadic grids
and grid stopping times. Each approximating policy has a strong
realization in the original model.
\end{lemma}
\begin{proof}
For mesh \(d_n\), define, on \([jd_n,(j+1)d_n)\), \(j\geq1\),
\[
 u_t^n=\frac1{d_n}\int_{(j-1)d_n}^{jd_n}u_s\,ds,
\]
and use any fixed element of \(U\) on the first interval. The
interval assumption on \(U\) makes these actions admissible and
known at the start of each interval. Translation continuity in
\(L^2(dt\,d\Qq)\), followed by Jensen's inequality, gives
\(u^n\to u\) in that space. Uniform continuity of \(\psi\) on \(U\)
gives \(\psi(u^n)\to\psi(u)\) in \(L^1\).
This is approximation of a given policy, not replacement of its
cost by the convex envelope of \(\psi\).

The states converge in \(L^2\) of the supremum norm; likelihoods
converge by boundedness of \(h\); quadratic terms are uniformly
integrable by Lemma~\ref{lem:Feller}. Round the bounded \(\tau\)
upward to the grid. It remains an observation stopping time,
and its deterministic bound grows by at most \(d_n\).
Continuity of paths and \(k\), with the same moment bounds,
gives convergence of the voluntary terminal term. The maturity
term converges by quadratic uniform integrability. The running
term converges by boundedness of \(\psi\), its convergence in
probability, and the likelihood moment bounds. Therefore
\[
 J_\pi(u^n,\tau^n)\longrightarrow J_\pi(u,\tau).
\]
If a terminal posterior function \(f\in\CW\) is used at a fixed
endpoint instead, its convergence follows from the conditional
measure argument of Lemma~\ref{lem:Feller}.

The history-dependent grid functional is strongly implementable
interval by interval as in Theorem~\ref{thm:grid}.
Its likelihood identifies its physical law with the reference
construction, without asserting a strong realization for the
arbitrary continuous-time reference policy.
\end{proof}

\begin{theorem}[The full value]\label{thm:value}
For every \(\pi\in\Ptwo\), the original value, the reference-policy
infimum, and the decreasing grid limit agree:
\begin{equation}\label{eq:Vlimit}
 V(\pi)=\lim_{n\to\infty}v_{2^{-n}}(\pi).
\end{equation}
The same is true for meshes \(t/2^n\), for any fixed \(t>0\).
For every initial posterior and every \(\varepsilon>0\), an exact
posterior policy on a sufficiently fine grid has original cost
at most \(V(\pi)+\varepsilon\).
\end{theorem}
\begin{proof}
The stopped representation of Lemma~\ref{lem:reverse} makes the
reference infimum no larger than the original value.
Every grid policy is original-admissible, so the original value
is no larger than any grid value. Conversely, truncate any
reference stopping rule using \eqref{eq:tail} and then use
Lemma~\ref{lem:density}. The infimum of grid values is no larger
than the cost of any reference policy. These three inequalities
give equality. Refinement enlarges the admissible grid policy
class, proving monotonicity. The density proof applies to either
nested mesh family, and the last assertion follows from
Theorem~\ref{thm:grid}.
\end{proof}

This is a convergent construction of pure policies; no uniform
numerical rate in the mesh is claimed under mere continuity of
\(k\). Each decision uses only the current posterior at that grid
date and the fixed clock, rather than an unobserved state or an
additional feature of the observation history.
These are stationary Markov decisions on each decision grid,
with the selected action held between checks. They are not asserted
to be continuous-time stationary feedbacks of the current
\(\pi_t\) alone between decision dates.

\section{Continuity and measurable continuation}

\begin{lemma}\label{lem:Vcontinuous}
The full value belongs to \(\CW\).
\end{lemma}
\begin{proof}
Immediate stopping gives the growth bound. Couple
\(\pi_n\to\pi\) in Wasserstein-\(2\) and use the same fixed reference
functional \(u(Y),\tau(Y)\) for both laws. On a fixed horizon,
\[
 \E_\Qq\sup_{s\leq T}|X_s^n-X_s|^2
       \leq C_T\E|X_0^n-X_0|^2
\]
holds uniformly over \(u\), hence also over bounded \(\tau\).
The bounded Lipschitz functions \(h,a\) give uniform convergence
in probability of the likelihood and survival factors.
On a compact state interval \(k\) is uniformly continuous; its
complement is controlled by the initial second-moment tails and
bounded-coefficient increments of Lemma~\ref{lem:Feller}, uniformly
over controls and stopping times. Conditional likelihood higher
moments are uniform too. Consequently
\[
 \sup_{u,\tau\leq T}|J_{\pi_n}(u,\tau)-J_\pi(u,\tau)|
             \longrightarrow0.
\]
The uniform tail bound \eqref{eq:tail} and boundedness of
\(\wt(\pi_n)\) extend this to all stopping times.
Taking infima proves continuity. No fourth initial moment is used.
\end{proof}

\begin{lemma}[Borel nearly optimal tails]\label{lem:paste}
For each \(\varepsilon>0\) there is a Borel family, indexed by the
restarting posterior \(\pi\), of strongly realizable grid policies
with costs at most \(V(\pi)+\varepsilon\wt(\pi)\).
\end{lemma}
\begin{proof}
Choose the first \(n\) such that
\(v_{2^{-n}}(\pi)\leq V(\pi)+\varepsilon\wt(\pi)\).
Continuity and \eqref{eq:Vlimit} make this index Borel.
For each of the countably many meshes use the fixed Borel selector
and sample update of Theorem~\ref{thm:grid}. At a restart, concatenate
that rule using fresh original Brownian increments.
\end{proof}

\begin{corollary}[The posterior immediate-stopping region]
\label{cor:stopregion}
The exact region where immediate voluntary stopping is optimal is
the closed set
\[
 D=\{\pi:V(\pi)=K(\pi)\}
   =\bigcap_{n\geq1}\{\pi:v_{2^{-n}}(\pi)=K(\pi)\}.
\]
The closed sets in this intersection decrease with \(n\).
\end{corollary}
\begin{proof}
We have \(V\leq v_{2^{-(n+1)}}\leq v_{2^{-n}}\leq K\).
Together with the pointwise limit, these inequalities prove the
identity and monotonicity. Continuity makes each set closed.
Stopping at once costs exactly \(K\), so it attains the value
exactly on \(D\).
\end{proof}

\section{The verified posterior obstacle}

\subsection{Greatest cost-submean minorant}

\begin{definition}
A weighted continuous cost-submean minorant is an \(f\in\CW\) such
that, for every \(\pi\in\Ptwo\), \(u\in U\), and \(t>0\),
\begin{equation}\label{eq:submean}
 f(\pi)\leq K(\pi),\qquad
 f(\pi)\leq g_t(\pi,u)+\cQ_t^u f(\pi).
\end{equation}
\end{definition}
No lower bound on \(f\) is required; its weighted norm is finite.

\begin{theorem}[Full measure-space obstacle]\label{thm:obstacle}
The value \(V\) is the greatest weighted continuous cost-submean
minorant. Thus \eqref{eq:submean}, with the greatest-element
requirement, has an existing and unique answer for every coefficient
set in the canonical problem.
\end{theorem}
\begin{proof}
We have \(V\leq K\). Keep constant control \(u\) to time \(t\) or
maturity, then concatenate the Borel policy of
Lemma~\ref{lem:paste}. The expected continuation excess is at most
\(\varepsilon e^{-\gamma t}\wt(\pi)\). Hence
\[
 V(\pi)\leq g_t(\pi,u)+\cQ_t^uV(\pi)
                  +\varepsilon e^{-\gamma t}\wt(\pi).
\]
Let \(\varepsilon\downarrow0\); continuity is
Lemma~\ref{lem:Vcontinuous}.

Conversely, \eqref{eq:submean} implies \(f\leq\cT_df\).
Monotonicity gives \(f\leq\cT_d^Nf\) for all \(N\); weighted
contraction then gives \(f\leq v_d\).
Let \(d=2^{-n}\downarrow0\) to obtain \(f\leq V\).
\end{proof}

This obstacle uses only known constant-control killed Bayes kernels
and \(K\). It does not define subsolutions through existence of an
optimal global strategy, and it does not assume a classical solution
or viscosity comparison theorem.

\subsection{A unique leavable-semigroup fixed point}

For \(f\in\CW\), \(0\leq f\leq K\), define
\begin{equation}\label{eq:H}
 \cH_tf=\lim_{m\to\infty}
           \big(\cT_{t/2^m}\big)^{2^m}f,\qquad t>0.
\end{equation}
The limit decreases pointwise under refinement.
Every finite composition is a finite-grid control/stopping problem
with terminal \(f\); thus the definition uses only explicit kernels.
If desired, its interpretation with arbitrary reference block rules
follows from Lemma~\ref{lem:density} and its terminal-posterior extension.
When a rounded stop reaches the endpoint, \(f\leq K\) prevents an
artificial increase in cost. This interpretation is not needed for
the fixed-point existence proof.

The finite compositions have contraction factor \(e^{-\gamma t}\).
Taking pointwise limits gives
\begin{equation}\label{eq:Hcontract}
 \|\cH_tf-\cH_tg\|_{\wt}
           \leq e^{-\gamma t}\|f-g\|_{\wt}.
\end{equation}

\begin{theorem}[Unique mild fixed point]\label{thm:mild}
Among \(F\in\CW\), \(0\leq F\leq K\), the value is the unique function
satisfying
\begin{equation}\label{eq:mild}
 F=\cH_tF\qquad(t>0).
\end{equation}
For each single \(t>0\), there is already at most one fixed point
in this class.
\end{theorem}
\begin{proof}
Theorem~\ref{thm:obstacle} gives \(V\leq\cT_dV\).
By monotonicity, the grid fixed-point identity, and \(V\leq v_d\),
\[
 V\leq\cT_d^NV\leq\cT_d^Nv_d=v_d.
\]
Take \(d=t/2^m\) and \(N=2^m\). Theorem~\ref{thm:value} gives
\(v_d\downarrow V\), so squeezing proves \(\cH_tV=V\).
This order argument and the strong measurable concatenation already
proved for the submean property supply existence without assuming
a continuous-time dynamic programming theorem.
If \(F,G\) are fixed points, \eqref{eq:Hcontract} gives
\(\|F-G\|_{\wt}\leq e^{-\gamma t}\|F-G\|_{\wt}\).
Their finite weighted distance must be zero.
\end{proof}

\section{Survival filtering and the differential obstacle}

For \(\varphi\in C_c^2(\R)\), the conditional Itô formula is
\begin{equation}\label{eq:Zakai}
 d\rho_t(\varphi)=
 \rho_t(\mathcal L^{u_t}\varphi-a\varphi)\,dt
                        +\rho_t(h\varphi)\,dY_t.
\end{equation}
The signal Brownian martingale has zero conditional projection
onto the observation filtration, and the observation-adapted
control factors out of the drift.
The quotient rule gives, on survival,
\begin{align}
 d\pi_t(\varphi)
 &=\big[\pi_t(\mathcal L^{u_t}\varphi)-\pi_t(a\varphi)
                   +\pi_t(a)\pi_t(\varphi)\big]\,dt\notag\\
 &\quad+
 \big[\pi_t(\varphi h)-\pi_t(\varphi)\pi_t(h)\big]\,dI_t,
 \qquad dI_t=dO_t-\pi_t(h)\,dt.                    \label{eq:KS}
\end{align}
One may retain \eqref{eq:Zakai} as the global reference construction.
No equality of observation and innovation filtrations is used.

For \(F(\pi)=f(\pi(\varphi_1),\ldots,\pi(\varphi_m))\) smooth
cylindrical, put
\begin{align*}
 b_i^u(\pi)&=\pi(\mathcal L^u\varphi_i)-\pi(a\varphi_i)
                                  +\pi(a)\pi(\varphi_i),\\
 \beta_i(\pi)&=\pi(\varphi_ih)-\pi(\varphi_i)\pi(h),\\
 \mathscr A^uF(\pi)&=\sum_i f_ib_i^u(\pi)
              +\frac12\sum_{i,j}f_{ij}\beta_i(\pi)\beta_j(\pi).
\end{align*}
The differential form of the verified mild obstacle is
\begin{equation}\label{eq:PDE}
 \min\left\{K(\pi)-V(\pi),\
 \inf_{u\in U}
 [\mathscr A^uV(\pi)-\pi(a)V(\pi)+\pi(ak)+\ell(u)]\right\}=0.
\end{equation}
For smooth cylinder tests, consistency follows from Itô's formula
and the kernels. Theorems~\ref{thm:obstacle} and \ref{thm:mild}
give the unconditional existence, verification, and uniqueness
in the stated mild senses. No uniqueness claim among unspecified
cylinder-viscosity solutions is made.

The drift term
\(-\pi(a\varphi)+\pi(a)\pi(\varphi)\) records survival information.
The killing source is \(\pi(ak)\), not \(\pi(a)\pi(k)\);
the latter would incorrectly omit the terminal state tilt.

\section{Exact-optimality certificate}

For a fixed reference control put
\begin{equation}\label{eq:M}
 M_t=\rho_t(1)V(\pi_t)+
          \int_0^t[\rho_s(ak)+\ell(u_s)\rho_s(1)]\,ds.
\end{equation}
This is a submartingale. Here is a derivation avoiding an assumed
continuous-time DPP. The submean inequality for \(V\), conditioned
at successive grid dates, implies the continuation inequality for
every observation-history grid control, with terminal \(V\).
For a prescribed progressive control approximate it by the
past-average grids in Lemma~\ref{lem:density}. Its terminal posterior
converges as in Lemma~\ref{lem:Feller}; continuity and the weight
bound for \(V\) permit passage to the limit. This proves the
continuation inequality for the prescribed control. Applying it
conditionally after each observed history, using the posterior flow,
gives the submartingale property. The known past is a fixed parameter;
the same measurable continuation construction applies.

Lemma~\ref{lem:versions} gives continuous posterior paths, hence a
continuous version of \eqref{eq:M}. The running integral is dominated
by its integrable infinite-time total. Also
\[
 \E_\Qq[\rho_t(1)V(\pi_t)]
       \leq C_ke^{-\gamma t}\wt(\pi)\longrightarrow0.
\]
Thus \(M_t\to M_\infty\) in \(L^1\), where \(M_\infty\) is the full
running integral in \eqref{eq:M}. Passing to this limit in the
submartingale inequality yields
\[
 0\leq M_t\leq\E_\Qq[M_\infty\mid\mathcal F_t^Y].
\]
It is consequently of class D, dominated by the closed martingale
of a fixed integrable variable. Alternatively, bounded-horizon
stopped terms have the uniform integrability of Section 4, and
their tails after \(T\) are controlled by the optional form of
\eqref{eq:mass}. Optional sampling is valid for unbounded voluntary
stopping times.

For every policy,
\begin{equation}\label{eq:certificate}
 \begin{split}
 J_\pi(u,\tau)-V(\pi)
  &=\E_\Qq[\rho_\tau(1)(K(\pi_\tau)-V(\pi_\tau))]\\
  &\quad+\E_\Qq[M_\tau-M_0],
 \end{split}
\end{equation}
with zero voluntary terminal term at infinity.
Both terms are nonnegative. Exact optimality is therefore equivalent
to contact at finite voluntary termination and zero expected
increase of the stopped Bellman submartingale. By class D, the
latter is equivalent to its martingale property up to termination.

In particular, an admissible posterior feedback making \(M\) a
martingale up to its first contact with \(\{V=K\}\), together with
that contact rule, is optimal. If additional Itô regularity of
\(V\) is justified, minimization of the continuation expression
in \eqref{eq:PDE} is a sufficient drift condition.
These are exact certificates, not an unsupported general assertion
of continuous-time pure attainment. The grid policies already
supply every \(\varepsilon\)-optimal rule under the original assumptions.

\section{A finite-dimensional structural class}

A verifiable sufficient structure within the original bounds is
as follows. Let \(h\) be constant. Suppose a finite-dimensional space
\(E\), spanned by \(q_1,\ldots,q_r\) satisfying
\(\int(1+x^2)|q_j(x)|\,dx<\infty\), is preserved by each
constant-control killed forward semigroup, with matrix
\(e^{tA(u)}\), where \(A(u)\) is continuous in \(u\).
Start with a nonnegative density in \(E\).
For a known piecewise constant control, successive matrix products
propagate the unnormalized survival density. Approximation of
bounded control paths gives
\[
 dZ_t=A(u_t)Z_t\,dt,\qquad
 \pi_t(x)=
 \frac{\sum_jZ_t^jq_j(x)}{\sum_jZ_t^j\int q_j}.
\]
The constant-\(h\) likelihood is scalar and cancels on normalization.
The actual killed semigroups preserve the nonnegative cone;
ODE convergence preserves it for general paths. State-law weak
convergence identifies the ODE limit with the true filter.
The control is known, so this is a finite-dimensional sufficient
posterior state.

With smooth bases the invariance can be checked by the killed
forward equation. For merely Lipschitz original coefficients,
adjoint notation is weak: the matrix identity means
\[
 \int q_j(\mathcal L^u\varphi-a\varphi)
          =\sum_i A_{ij}(u)\int q_i\varphi
       \quad(\varphi\in C_c^2),
\]
together with the forward-semigroup invariance, not an unjustified
pointwise derivative of \(\sigma\). The next example identifies
its actual killed density directly and requires no general SPDE
uniqueness theorem.

This is a sufficient structural subclass only. The general theorem
retains arbitrary original \(h,a\). In particular the usual
nonconstant linear Kalman observation would violate boundedness
of \(h\) here and is not invoked.

\section{An exact nonconstant-hazard filter and stopping region}

Take
\[
 \mu=0,\quad \sigma=1,\quad h=0,\quad
 a(x)=7-6\sech^2x\geq1,
\]
and
\[
 q_0(x)=\sech^3x,\qquad
 q_2(x)=\sech x(5\tanh^2x-1)=4\sech x-5\sech^3x.
\]
Write
\[
 I_0=\int q_0=\pi/2,\quad I_2=\int q_2=3\pi/2,\quad
 p_z(x)=\frac{q_0(x)+zq_2(x)}{I_0+zI_2},\qquad0\leq z\leq1.
\]
The \(\pi\) in these constants is the circle constant.
Since \(q_2/q_0=4\cosh^2x-5\geq-1\), these are probability densities
with finite second moments. All coefficient assumptions hold.
Direct differentiation gives
\[
 \tfrac12q_0''-a q_0=-\tfrac52q_0,\qquad
 \tfrac12q_2''-a q_2=-\tfrac{13}{2}q_2.
\]
The unnormalized survival density is consequently
\[
 \frac{e^{-5t/2}q_0+ze^{-13t/2}q_2}{I_0+zI_2}.
\]
This is the actual killed density, not only an invariant ansatz:
it solves the killed heat equation with the prescribed initial law.
The Duhamel equation with the heat semigroup is unique in \(L^1\),
since the norm of a difference is bounded by \(\|a\|_\infty\)
times the integral of that norm. Gronwall's inequality makes it zero.
The Feynman--Kac density solves the same equation.

Thus the posterior on survival is \(p_{z_t}\), with
\[
 z_t=ze^{-4t},\qquad
 p_z(a)=\frac{5+39z}{2(1+3z)}.
\]
For \(z>0\), survival changes this posterior strictly.
The Brownian observation is uninformative but state-dependent
maturity is informative.

Here the state drift is control-independent, so choose any minimizer
of \(\psi\). Put
\[
 \ell_*=\min_{u\in U}\psi(u)+c>0,\qquad
 k(x)=4\ell_*\sech^2x.
\]
The immediate-stop cost is
\[
 K(z)=\ell_*\frac{3+z}{1+3z}.
\]
For general continuous quadratic-growth \(k\), define
\[
 K_j=\int kq_j,\quad A_j=\int akq_j,\quad
 \lambda_0=5/2,\quad\lambda_2=13/2.
\]
Integrating the two killed eigenmodes gives the cost at deterministic
voluntary time \(t\):
\begin{equation}\label{eq:Fgeneral}
 \begin{split}
 F_z(t)=\frac1{I_0+zI_2}\bigg[
 &e^{-\lambda_0t}K_0+ze^{-\lambda_2t}K_2\\
 &+\frac{A_0+\ell_*I_0}{\lambda_0}(1-e^{-\lambda_0t})\\
 &+z\frac{A_2+\ell_*I_2}{\lambda_2}(1-e^{-\lambda_2t})
 \bigg].
 \end{split}
\end{equation}
For the displayed \(k\), elementary \(\sech\)-power integrals,
or integration by parts, give
\[
 B_0:=K_0-\frac{A_0+\ell_*I_0}{\lambda_0}
               =\frac{\pi\ell_*}{10},\qquad
 B_2:=K_2-\frac{A_2+\ell_*I_2}{\lambda_2}
               =-\frac{\pi\ell_*}{2}.
\]
Thus
\begin{equation}\label{eq:Fexplicit}
 F_z(t)=\frac{\ell_*}{1+3z}
 \left[\frac{14}{5}+2z+\frac15e^{-5t/2}
                                -ze^{-13t/2}\right].
\end{equation}
Its derivative has the sign of \(13ze^{-4t}-1\).
It has at most one finite critical point, and when present that
point is a maximum. Its minimum is at \(0\) or infinity.

A pre-default observation stop here is a functional of an independent
Brownian path. Its cost is an average of these deterministic-time
costs, so randomization cannot improve the minimum. The exact answer is
\begin{equation}\label{eq:example}
 V(z)=\frac{\ell_*}{1+3z}
          \min\left\{3+z,\frac{14}{5}+2z\right\}.
\end{equation}
Stop immediately for \(z\geq1/5\); for \(z<1/5\), never stop
voluntarily and let maturity terminate the problem.
At \(z=1/5\) either endpoint rule is optimal. Since \(z_t\) decreases,
the continuation rule never enters the stopping region. Maturity
has mean at most \(1\). This is a valid changing survival filter
and an explicit posterior stopping threshold within the original
bounded-coefficient model.

\section{Full-scope verification and provenance}

The general result quantifies over every prior in \(\Ptwo\), every
compact control interval, every original bounded Lipschitz
coefficient set, and every stated continuous cost. It retains the
infinite horizon and state-dependent random maturity. Controls
and stopping rules use only observations and survival.
The canonical requirements are met as follows.
\begin{enumerate}
\item The full value is Theorem~\ref{thm:value}: the original
admissible infimum, not merely a relaxed or independent-hazard value.
\item Theorem~\ref{thm:grid} gives posterior decision and stopping
rules attaining all grid values, and Theorem~\ref{thm:value} supplies
every requested \(\varepsilon\)-optimal original rule.
Equation~\eqref{eq:certificate} characterizes exact optimality
when attained.
\item Survival information and terminal maturity tilt are
\eqref{eq:KS} and \eqref{eq:tilt}.
\item The general measure-valued obstacle is solved and verified
by Theorems~\ref{thm:obstacle} and \ref{thm:mild};
\eqref{eq:PDE} is its differential form.
Neither an unproved classical-solution assumption nor a viscosity
comparison assertion replaces that proof.
\item Section 11 supplies finite-dimensional structural assumptions,
and Section 12 gives a state-dependent-hazard filter with an exact
posterior stopping region.
\end{enumerate}

The agenda source is Beneš, Gaitsgori and Karatzas \cite{BGK},
Section 5. The combined concrete model is the canonical C61-73 model,
not a theorem already in that source. Filtering identities are
contextualized by Fujisaki, Kallianpur and Kunita \cite{FKK}; the killed
filter and maturity tilt have been derived here. Bandini, Cosso,
Fuhrman and Pham \cite{BCFP} discuss posterior Bellman equations and
the Nisio semigroup/mild-solution approach. Their viscosity comparison
theorem is not imported into this killed stopping problem.

This hand-proof manuscript was prepared with GPT Codex assistance.
No Lean formalization or external peer-review certification is claimed.

\begin{thebibliography}{9}
\bibitem{BGK}
V. E. Beneš, G. Gaitsgori and I. Karatzas,
\emph{Drift Control with Discretionary Stopping for a Diffusion},
arXiv:2401.10043v3, Section 5.
\url{https://arxiv.org/html/2401.10043v3#S5}.

\bibitem{FKK}
M. Fujisaki, G. Kallianpur and H. Kunita,
\emph{Stochastic differential equations for the non linear filtering
problem}, Osaka Journal of Mathematics \textbf{9} (1972), 19--40.
\url{https://doi.org/10.18910/5728}.

\bibitem{BCFP}
E. Bandini, A. Cosso, M. Fuhrman and H. Pham,
\emph{Randomized filtering and Bellman equation in Wasserstein space
for partial observation control problem}, arXiv:1609.02697;
Stochastic Processes and their Applications \textbf{129} (2019), 674--711.
\url{https://arxiv.org/html/1609.02697}.
\url{https://doi.org/10.1016/j.spa.2018.03.014}.
\end{thebibliography}
\end{document}
